\begin{ExerciseBlock}
\ExerciseHeading{02}{Schwerpunkt einer Lochplatte}{rigid_bodies_perforated_plate}
Eine rechteckige Platte (Länge $a = 0,6\,$m, Breite $b = 0,3\,$m) hat im Abstand $c = 0,142\,$m von der Mitte auf einer Seite ein Loch mit Radius $r = 0,1\,$m. Die Masse dieser Lochplatte beträgt $M = 47,3\,$kg.

\begin{enumerate}
	\item Berechnen Sie die Masse einer Scheibe,
die bei gleicher Dicke genau in das Loch
passen würde.
	\item Berechnen Sie die Koordinate $x_s$ des Schwerpunktes im angegebenen Koordinatensystem.
	\item Stellen Sie eine Formel für das Trägheitsmoment $J_z$ dieses Bauteils für Rotation um die $z$-Achse auf.

\end{enumerate}

\begin{center}
	\CachedTikzFigure{content/resources/chapters/rigid_bodies/tikz/perforated_plate.tex}
\end{center}
\end{ExerciseBlock}
\UseExerciseSolution{rigid_bodies_perforated_plate}
